%Version 3.1 December 2024
% See section 11 of the User Manual for version history
%
%%%%%%%%%%%%%%%%%%%%%%%%%%%%%%%%%%%%%%%%%%%%%%%%%%%%%%%%%%%%%%%%%%%%%%
%%                                                                 %%
%% Please do not use \input{...} to include other tex files.       %%
%% Submit your LaTeX manuscript as one .tex document.              %%
%%                                                                 %%
%% All additional figures and files should be attached             %%
%% separately and not embedded in the \TeX\ document itself.       %%
%%                                                                 %%
%%%%%%%%%%%%%%%%%%%%%%%%%%%%%%%%%%%%%%%%%%%%%%%%%%%%%%%%%%%%%%%%%%%%%

%%\documentclass[referee,sn-basic]{sn-jnl}% referee option is meant for double line spacing

%%=======================================================%%
%% to print line numbers in the margin use lineno option %%
%%=======================================================%%

%%\documentclass[lineno,pdflatex,sn-basic]{sn-jnl}% Basic Springer Nature Reference Style/Chemistry Reference Style

%%=========================================================================================%%
%% the documentclass is set to pdflatex as default. You can delete it if not appropriate.  %%
%%=========================================================================================%%

% \documentclass[sn-basic]{sn-jnl}% Basic Springer Nature Reference Style/Chemistry Reference Style

%%Note: the following reference styles support Namedate and Numbered referencing. By default the style follows the most common style. To switch between the options you can add or remove “Numbered” in the optional parenthesis. 
%%The option is available for: sn-basic.bst, sn-chicago.bst%  
 
%%\documentclass[pdflatex,sn-nature]{sn-jnl}% Style for submissions to Nature Portfolio journals
% \documentclass[pdflatex,sn-basic]{sn-jnl}% Basic Springer Nature Reference Style/Chemistry Reference Style
\documentclass[pdflatex,sn-mathphys-num]{sn-jnl}% Math and Physical Sciences Numbered Reference Style
% Source-review markers are no-op by default and are not printed.
\providecommand{\reviewermap}[3]{}
%%\documentclass[pdflatex,sn-mathphys-ay]{sn-jnl}% Math and Physical Sciences Author Year Reference Style
%%\documentclass[pdflatex,sn-aps]{sn-jnl}% American Physical Society (APS) Reference Style
%%\documentclass[pdflatex,sn-vancouver-num]{sn-jnl}% Vancouver Numbered Reference Style
%%\documentclass[pdflatex,sn-vancouver-ay]{sn-jnl}% Vancouver Author Year Reference Style
%%\documentclass[pdflatex,sn-apa]{sn-jnl}% APA Reference Style
%%\documentclass[pdflatex,sn-chicago]{sn-jnl}% Chicago-based Humanities Reference Style

%%%% Standard Packages
%%<additional latex packages if required can be included here>

\usepackage{graphicx}%
\usepackage{multirow}%
\usepackage{amsmath,amssymb,amsfonts}%
\usepackage{amsthm}%
\usepackage{mathrsfs}%
\usepackage[title]{appendix}%
\usepackage{xcolor}%
\usepackage{textcomp}%
\usepackage{manyfoot}%
\usepackage{booktabs}%
\usepackage{algorithm}%
\usepackage{algorithmicx}%
\usepackage{algpseudocode}%
\usepackage{listings}%

% User defined
%%%%Userd-added package
\usepackage[version=3]{mhchem}
\usepackage{subcaption}
\usepackage{graphicx}
\usepackage{float}
\usepackage{multirow}
\usepackage{array}
\usepackage{amsmath}
\usepackage{booktabs}
\usepackage{amsmath}
\usepackage{amssymb}
\usepackage{xcolor}

\usepackage[utf8]{inputenc} % allow utf-8 input
\usepackage{hyperref}       % hyperlinks
\usepackage{url}            % simple URL typesetting
\usepackage{amsfonts}       % blackboard math symbols
\usepackage{nicefrac}       % compact symbols for 1/2, etc.
\usepackage{microtype}      % microtypography

\usepackage{xr}
\externaldocument{sn-article-SI}
%%%%
\newcommand{\pc}[1]{{\color{purple}[Paulette: #1]}}
\newcommand{\yi}[1]{{\color{teal}[Yi: #1]}}
%%%%%=============================================================================%%%%
%%%%  Remarks: This template is provided to aid authors with the preparation
%%%%  of original research articles intended for submission to journals published 
%%%%  by Springer Nature. The guidance has been prepared in partnership with 
%%%%  production teams to conform to Springer Nature technical requirements. 
%%%%  Editorial and presentation requirements differ among journal portfolios and 
%%%%  research disciplines. You may find sections in this template are irrelevant 
%%%%  to your work and are empowered to omit any such section if allowed by the 
%%%%  journal you intend to submit to. The submission guidelines and policies 
%%%%  of the journal take precedence. A detailed User Manual is available in the 
%%%%  template package for technical guidance.
%%%%%=============================================================================%%%%

%% as per the requirement new theorem styles can be included as shown below
\theoremstyle{thmstyleone}%
\newtheorem{theorem}{Theorem}%  meant for continuous numbers
%%\newtheorem{theorem}{Theorem}[section]% meant for sectionwise numbers
%% optional argument [theorem] produces theorem numbering sequence instead of independent numbers for Proposition
\newtheorem{proposition}[theorem]{Proposition}% 
%%\newtheorem{proposition}{Proposition}% to get separate numbers for theorem and proposition etc.

\theoremstyle{thmstyletwo}%
\newtheorem{example}{Example}%
\newtheorem{remark}{Remark}%

\theoremstyle{thmstylethree}%
\newtheorem{definition}{Definition}%

\raggedbottom
%%\unnumbered% uncomment this for unnumbered level heads

\begin{document}

%% Previous working title removed for REV-A4 scope control.
\title[Article Title]{Non-Equilibrium Properties Reveal the Limits of Specialist and Foundation Machine-Learned Force Fields}
%%=============================================================%%
%% GivenName	-> \fnm{Joergen W.}
%% Particle	-> \spfx{van der} -> surname prefix
%% FamilyName	-> \sur{Ploeg}
%% Suffix	-> \sfx{IV}
%% \author*[1,2]{\fnm{Joergen W.} \spfx{van der} \sur{Ploeg} 
%%  \sfx{IV}}\email{iauthor@gmail.com}
%%=============================================================%%

\author[1]{\fnm{Yi} \sur{Cao}}\email{ycao73@jh.edu}

\author[1]{\fnm{Peter} \sur{Mastracco}}\email{pmastra2@jh.edu}
% \author[1]{\fnm{Victor} \sur{Wu}}\email{vwu16@jhu.edu}
% % \equalcont{These authors contributed equally to this work.}

\author*[1]{\fnm{Paulette} \sur{Clancy}}\email{pclancy3@jhu.edu}
% \equalcont{These authors contributed equally to this work.}

\affil[1]{\orgdiv{Department of Chemical and Biomolecular Engineering}, \orgname{ Johns Hopkins University}, \orgaddress{\street{3400 N. Charles Street}, \city{Baltimore}, \postcode{21218}, \state{Maryland}, \country{USA}}}

%%==================================%%
%% Sample for unstructured abstract %%
%%==================================%%

% \abstract{The abstract serves both as a general introduction to the topic and as a brief, non-technical summary of the main results and their implications. Authors are advised to check the author instructions for the journal they are submitting to for word limits and if structural elements like subheadings, citations, or equations are permitted.}

\abstract{Machine-learned force fields (MLFFs), particularly those arising from pre-trained foundation models, offer an enticing compromise between \textit{ab initio} accuracy and the larger scope possible with semi-empirical force field-derived  molecular dynamics with respect to length and time scales. Despite their promise, a fundamental question remains: To what extent should one rely on specialist models, fine-tuned foundation models, or a hybrid approach? Trade-offs in data efficiency, robustness, and non-equilibrium fidelity remain poorly understood, as does understanding how different training paradigms encode physical information.
To assist with this decision, we introduce a unified benchmarking framework based on non-equilibrium configurational probes that span interpolation and extrapolation regimes. Migration pathways computed via Nudged Elastic Band calculations serve as the primary diagnostic probe, complemented by equilibrium and mechanical tests. Using Cr-doped \ce{Sb2Te3} as a representative 2D material for which non-equilibrium properties are critically important, we evaluate multiple training strategies within the MACE architecture.
While all models reproduce equilibrium structures, they diverge sharply in their ability to adequately capture non-equilibrium processes. Fine-tuned MACE models improve selected kinetic probes but can degrade long-range mechanical response. ``Zero-shot'' MACE foundation models exhibit overly smoothed potentials and specialist models can fail for high-energy transitions. Latent-space and surrogate-error diagnostics are consistent with distinct model representations across training paradigms, supporting the conclusion that equilibrium metrics alone are insufficient.

\reviewermap{REV-A4}{done-local}{Scope generalizable-standard language to a candidate framework demonstrated on Cr-doped Sb2Te3.}
This work demonstrates migration-based non-equilibrium probes as a data-efficient candidate framework for MLFF evaluation, developed and tested on a representative material system for which such probes are important, and provides practical guidance for robust MLFF development.
}
%%================================%%
%% Sample for structured abstract %%
%%================================%%

% \abstract{\textbf{Purpose:} The abstract serves both as a general introduction to the topic and as a brief, non-technical summary of the main results and their implications. The abstract must not include subheadings (unless expressly permitted in the journal's Instructions to Authors), equations or citations. As a guide the abstract should not exceed 200 words. Most journals do not set a hard limit however authors are advised to check the author instructions for the journal they are submitting to.
% 
% \textbf{Methods:} The abstract serves both as a general introduction to the topic and as a brief, non-technical summary of the main results and their implications. The abstract must not include subheadings (unless expressly permitted in the journal's Instructions to Authors), equations or citations. As a guide the abstract should not exceed 200 words. Most journals do not set a hard limit however authors are advised to check the author instructions for the journal they are submitting to.
% 
% \textbf{Results:} The abstract serves both as a general introduction to the topic and as a brief, non-technical summary of the main results and their implications. The abstract must not include subheadings (unless expressly permitted in the journal's Instructions to Authors), equations or citations. As a guide the abstract should not exceed 200 words. Most journals do not set a hard limit however authors are advised to check the author instructions for the journal they are submitting to.
% 
% \textbf{Conclusion:} The abstract serves both as a general introduction to the topic and as a brief, non-technical summary of the main results and their implications. The abstract must not include subheadings (unless expressly permitted in the journal's Instructions to Authors), equations or citations. As a guide the abstract should not exceed 200 words. Most journals do not set a hard limit however authors are advised to check the author instructions for the journal they are submitting to.}

\keywords{Machine-learned force fields (MLFFs), foundation models, fine-tuning, defect migration pathways, 2D materials}

%%\pacs[JEL Classification]{D8, H51}

%%\pacs[MSC Classification]{35A01, 65L10, 65L12, 65L20, 65L70}

\maketitle

\section{Introduction}

Machine-Learned Force Fields (MLFFs)~\citep{unke_machine_2021, ssmith_ani-1_2017, chen_accurate_2017, batatia_mace_nodate,vandermause_active_2022, batzner_e3-equivariant_2022, bartok_gaussian_2010, friederich2021machine, zeng_deepmd-kit_2023, wines_chips-ff_2025, musaelian_learning_2023, chen_universal_2022, lindsey2017chimes,  xie2018crystal} have emerged as a transformative solution to the dilemma of accuracy-versus-cost in molecular simulation. The development of large-scale, pre-trained so-called ``foundation models,'' such as CHGNet,~\citep{deng_chgnet_2023} MEGNet,~\citep{chen_graph_2019} M3GNet,~\citep{chen_universal_2022} MACE-MP,~\citep{batatia_mace_nodate} MatterSim,~\citep{yang_mattersim_2024} and UMA,~\citep{wood_uma_2025} marks a paradigm shift to try to resolve this dilemma.~\citep{merchant_scaling_2023}.
However, this shift presents researchers with a critical question: For a specific system, is it more effective to invest significant resources to build a robust \textit{specialist} model from scratch,~\citep{zhang_deep_2018, smith_approaching_2017, chmiela_sgdml_2017} or to adapt a \textit{generalist} foundation model through fine-tuning? While generic benchmarks like Matbench,~\citep{dunn_benchmarking_2020} JARVIS-Leaderboard,~\citep{choudhary2024jarvis} and LeMat-GenBench,~\citep{betala2025lemat} validate their broad utility to predict equilibrium properties for which they were trained, 
their reliability for specific, high-fidelity applications remains an open question.~\citep{friederich_machine-learned_2021, deringer_machine_2019, jacobs2025practical}  
A valid concern is that these models fail to capture potentially more nuanced interactions that govern critical processes that shape the materials and their processing. ~\citep{schran_machine_2020, kapil_assessment_2022, grisafi_transferable_2018, liu2023discrepancies} 
Fine-tuning is a promising solution, but using naïve strategies can lead to \textit{``catastrophic forgetting''}, where general knowledge about the system is erased. Key questions of data efficiency,~\citep{janet_designing_2019, lookman_active_2019} stability,~\citep{zuo_performance_2020, jinnouchi_phase_2020} and transferability~\citep{zeni_data_2021, vandermause_on--fly_2020} for non-equilibrium processes remain largely unsolved. ~\citep{bernstein_quantifying_2019, zhang_active_2019, he2018statistical}



% REV-FORMALIZATION-2026-09-12 BEGIN: replace vague benchmark motivation with measurement-distribution framing.
% Original paragraph began: "To build a rigorous benchmark that addresses the specialist-generalist dilemma..."
To build a rigorous benchmark that addresses the dilemma facing the specialist versus the generalist, we must first define what distribution and what observables are being measured. Conventional MLFF validation typically estimates pointwise energy and force errors on near-equilibrium configurations. Such tests are necessary, but they do not, by themselves, measure whether the learned potential-energy surface is reliable along the low-probability configurations that control kinetic and mechanical processes.
% REV-FORMALIZATION-2026-09-12 END.

%%%%% updated for broader tone discussing about the non-equilibrium data's importance 
% REV-FORMALIZATION-2026-09-12 BEGIN: define pathways as task-conditioned probes and point readers to formal section.
% Original paragraph began: "We propose that overcoming this limitation requires evaluation strategies..."
To test this hypothesis, we formulate migration and collective-displacement pathways as task-conditioned [Yi- what is meant by a "task-conditioned probe"? The wording is rather non-specific] probes of the potential energy surface. In this view, a pathway is not merely an application-specific calculation; it is a structured evaluation distribution that explicitly samples configurations between metastable states, including transition-state and strained-lattice regions that may well be rare or absent in equilibrium datasets. [Yi- That sentence is far too long and it is not direct enough for me to understand what you intend to say. Let's  discuss. ] The mathematical framework used to separate equilibrium pointwise accuracy, fixed-path energetic fidelity, and self-relaxed path stability is defined in Sec.~\ref{sec:benchmark_formalization}.
% REV-FORMALIZATION-2026-09-12 END.

To demonstrate and validate this benchmarking framework, we employ doped two-dimensional (2D) van der Waals (vdW) materials as an ideal test bed. The discovery and design of novel 2D vdW materials~\citep{novoselov_2d_2016, manzeli_2d_2017, akinwande2017review, guo_stacking_2021, liu_roadmap_2024} continues to drive innovation in spintronics,~\citep{he_topological_2022} quantum devices,~\citep{qian_van_2024, song_electron_2018, nayak_microsoft_2025} and energy technologies.~\citep{jin_flexible_2019,gautam_creation_2024} A particularly powerful strategy for tuning the properties of these layered materials involves doping, the insertion of ``foreign'' atoms between vdW layers.~\citep{novoselov_2d_2016, whittingham_lithium_2004} This approach has enabled remarkable advances in materials design, ranging from lithium-ion batteries~\citep{winter_insertion_1998, hu_lithium_2024} to the engineering of magnetic semiconductors~\citep{gibertini_magnetic_2019} and quantum anomalous Hall systems.~\citep{tokura_magnetic_2019} Crucially, their layered structure offers distinct environments for atomic motion, making them an excellent probe for testing MLFF capabilities across different chemical contexts within a single system.

Among 2D materials, topological insulators such as antimony telluride (\ce{Sb2Te3})~\citep{zhang_topological_2009, guo_roadmap_2015} represent an especially promising class for doping studies. The doping of transition metals has proven to be effective across multiple topological systems: Iron in \ce{Bi2Se3} creates ferromagnetic order~\citep{checkelsky_trajectory_2014, kulbachinskii_thermoelectric_2012, haazen_ferromagnetism_2012}. Manganese in \ce{Bi2Te3} enables tunable magnetic properties.~\citep{klimovskikh_tunable_2020, fu_electronic_2021} And vanadium in \ce{Sb2Te3} modifies its electronic structure.~\citep{sun_sb2te3_2023, zhang_quantum_2021} The doping of chromium (Cr) into \ce{Sb2Te3}, the focus of this work, offers a compelling route to engineer its magnetic and topological properties for spintronic applications.~\citep{cortie_creating_2020,han_effect_2013,de_a_deus_magnetic_2021} Predicting the stability and dynamics of these complex systems is essential for guiding experimental synthesis, a task for which large-scale Molecular Dynamics (MD) simulations are extremely powerful~\citep{Zhang2023TuningStudy,thompson_lammps_2022} provided that a suitably representative force field is available. However, the atomic-scale processes governing functionality, including dopant migration and thermal stability, occur on time- and length- scales far beyond the reach of first-principles methods like density functional theory (DFT).~\citep{kohn_self-consistent_1965, perdew_generalized_1996, sohier_density_2017,choudhary_high-throughput_2017}

This work addresses these gaps by presenting a systematic benchmark of bespoke training and fine-tuning strategies for the case study of Cr-doped \ce{Sb2Te3} using the MACE architecture (Fig.~\ref{fig:conceptual_diagram}). By combining latent-space analysis with migration-based probes, we move beyond traditional accuracy metrics to evaluate the dynamic stability and kinetic predictions of each model. Our study aims to serve as both a practical guide for selecting MLFF training strategies and a diagnostic framework for designing more data-efficient learning loops. 
\reviewermap{REV-A7}{done-local}{Add nearest-neighbor literature and state the delta relative to the earlier arXiv/workshop report.}
This manuscript extends our earlier report~\citep{cao2025migration} by adding a unified fixed-path and self-consistent NEB analysis, interlayer sliding tests, representation diagnostics, and a data provenance workflow. It therefore treats migration pathways not only as an application result, but also as a structured benchmark distribution for testing model behavior under task-relevant distribution shift. [Yi- I think this sentence was repeated earlier. It needs rewording to be easier to understand.] This revised framing also positions the work relative to recent studies of MACE fine-tuning for ion diffusion, foundation-MLIP migration-barrier evaluation, and systematic softening of universal MLIP potential-energy surfaces [PLACEHOLDER: final nearest-neighbor literature citation-key audit].[Yi- needs a reference and a bit more explanation: How is this work similar/different to other studies?]

To this end, we address the following key questions:
\begin{enumerate}
\item How do bespoke MLFFs, trained from scratch, compare against large pre-trained foundation models in terms of accuracy, stability, and data efficiency for molecular simulation? What are the performance and stability trade-offs of different fine-tuning strategies?
\reviewermap{REV-A4}{done-local}{Frame generalizability as an open question rather than a completed claim.}
\item Can non-equilibrium probes generated from methods, such as NEB, provide an efficient benchmark for specialist versus generalist models, and what additional cross-system evidence is needed before such probes can be considered generalizable across materials and architectures?
\item Can latent-space analysis reveal failure signatures to guide data acquisition?
\end{enumerate}



\begin{figure}[htbp]
    \centering
    \includegraphics[width=0.7\textwidth]{Fig/NeurIPS2025-AI4MAT-fig1-schematic.pdf}
    \caption{\textbf{Conceptual Workflow Diagram} illustrating two competing workflows: (a) training a specialist MLFF from scratch and (b) fine-tuning a generalist foundation model. The diagram highlights the benchmarking questions addressed in this paper.}
    \label{fig:conceptual_diagram}
\end{figure}



\section{Methodology}

We employed a systematic approach to benchmark various Machine-Learned Force Fields (MLFFs) for Cr-doped \ce{Sb2Te3}\footnote{This material system serves as an ideal testbed because its anisotropic van der Waals structure provides distinct migration environments: in-gap diffusion within quintuple layers (representing in-distribution scenarios) and interlayer migration across vdW gaps (probing out-of-distribution regimes). This structural duality enables efficient sampling of diverse configurations—from equilibrium states to high-energy transition states—while addressing technologically relevant processes where both thermodynamic stability and kinetic properties are critical for doping engineering.}, 
comparing specialist models trained from scratch against fine-tuned foundation models.


\subsection{Data-Generation and Model-Training}

Reference data were generated using density functional theory (DFT) calculations adopting the well-used PBE functional~\citep{perdew_generalized_1996} and comprising roughly 20,000 configurations of a $4\times 2 \times 1$ system of the \ce{Sb2Te3} crystal~\citep{jain2013commentary}. We use a  1~fs timestep  
for a series of \textit{ab initio} molecular dynamics (AIMD) simulations at three temperatures (300~K, 600~K, 1200~K). We employed the MACE architecture for all models, which is a message-passing neural network for the representation of atomic interactions.~\cite{batatia_mace_nodate} 


This paper evaluates four distinct training strategies to benchmark their performance, named as follows:
\begin{enumerate}
\item \textbf{Scratch}: MACE model trained exclusively on our AIMD dataset
\item \textbf{Foundation}: Pre-trained MACE-MP model used without fine-tuning
\item \textbf{FT-600~K}: Foundation model fine-tuned on 5\% of data from 600~K trajectories
\item \textbf{FT-MultiT}: Foundation model fine-tuned on multi-temperature data from the three temperatures studied
\end{enumerate}

\reviewermap{REV-A2/REV-C10}{done-local}{Grouped split, train--test proximity audit, and 3-seed retraining are complete for FT-600K (OMAT-FT and MPA0 starts), Scratch, and Scratch-5\%; updated RMSE and in-gap barrier statistics are reported below and in Table~S1.}
Revision~1 treats frame-level train/test splitting as insufficient for benchmark claims involving migration pathways. The revised data workflow assigns each entire source trajectory to a single partition (train, validation, or test) rather than shuffling individual frames, so that no trajectory is split across partitions. Under this grouped protocol (1,786/311/267 train/validation/test frames), the nearest-neighbor SOAP cosine distance between any in-gap NEB image and the grouped training set is $5.4\times10^{-4}$, and the closest comparable Cartesian RMSD is 2.06~\AA; both are far outside the exact/near-duplicate thresholds used for this audit (SOAP $\le 10^{-4}$, RMSD $\le 0.05$~\AA), so no NEB-image/training-frame overlap was detected under the declared metrics. The Scratch-5\% training subset (89 frames) gives the same result (SOAP $6.4\times10^{-4}$, RMSD 2.06~\AA). Grouped-split retraining across seeds 123/234/345 is complete for FT-600K fine-tuned from the OMAT-FT checkpoint, Scratch, and Scratch-5\%, together with a single-seed MACE-MPA-0 fine-tuned comparator (excluded from 3-seed uncertainty estimates); the resulting RMSE and in-gap barrier statistics are reported below and in Table~S1.

\reviewermap{REV-C8}{done-local}{Zero-force NEB exporters are excluded from reported training provenance; release hash \texttt{f5c6d19} recorded.}
The reported fine-tuned models are documented as being trained on AIMD frames with real DFT energies and forces, not on NEB-derived exports lacking real force labels. The revised data manifest records source trajectory, split assignment, energy/force key mapping, and checksums for model-training files. The provenance audit covers 9,163 frames across the historical and grouped exports; the only nine exact zero-force frames are explicitly labelled one-atom Cr/Sb/Te isolated-atom $E_0$ references, with zero unclassified and zero NEB-derived zero-force frames. The complete manifest and audit code are in release commit \texttt{f5c6d19}.

\subsection{Evaluation Framework}
% REV-FORMALIZATION-2026-09-12 BEGIN: insert measurement-operator framework and external formalization module.
% Original paragraph began: "Models were assessed through: (i) Molecular dynamics simulations..."
Models were assessed through a hierarchy of measurement operators that separate near-equilibrium pointwise accuracy from task-conditioned process fidelity. Static energy and force errors, molecular dynamics observables, fixed-path NEB scoring, self-relaxed NEB stability, interlayer sliding, and representation diagnostics were evaluated as distinct probes of the same candidate force fields. Detailed computational parameters are provided in the Supplementary Information.

\input{revision1/benchmark_formalization}
% REV-FORMALIZATION-2026-09-12 END.

\subsection{SHAP Analysis}

\reviewermap{REV-A6/REV-C13}{done-local}{SHAP methods updated to the audited 150-frame absolute-energy-error surrogate workflow; X-FORCE three-model source lineage and release evidence are documented in \texttt{docs/audit/shap\_source\_lineage.md}.}

We employed a model-agnostic SHAP workflow to identify SOAP descriptors associated with MACE energy-error behavior. Local atomic environments were encoded using SOAP descriptors (5.0~\AA\ cutoff, $n_{\mathrm{max}}=4$, $l_{\mathrm{max}}=4$), capturing species-resolved radial and angular interactions in Cr-doped \ce{Sb2Te3}. For each MACE variant, gradient-boosted regressors (200 trees, learning rate 0.05, max depth 8) were trained to predict the absolute total-energy error on the audited 150-frame trajectory. TreeSHAP values were then computed for the fitted surrogate models, and surrogate fidelity was assessed with 5-fold cross-validation rather than in-sample $R^2$. We compared the top 50 influential SOAP features across models to identify shared versus model-specific sensitivities in the surrogate error model.


\section{Results and Discussions}

\subsection{Equilibrium Properties: Foundation for Comparison}
% REV-FORMALIZATION-2026-09-12 BEGIN: clarify that equilibrium RMSE estimates pointwise error under mu_eq.
% Original paragraph began: "We considered four candidate models..."
As a baseline, we first evaluated all four candidate models on the near-equilibrium DFT data used in conventional MLFF validation. In the notation of Sec.~\ref{sec:benchmark_formalization}, this corresponds to estimating a pointwise error under $\mu_{\mathrm{eq}}$, the distribution of relaxed and thermally sampled configurations near local minima. All models satisfy typical energy and force RMSE thresholds (Table~S1), and the corresponding parity plots show visually similar agreement with DFT across energies and forces (Fig.~S1).

The convergence of these pointwise metrics does not imply equivalent behavior under the task-conditioned distributions $\mu_\tau$ associated with migration, sliding, or rare-event kinetics. We therefore treat equilibrium accuracy as a necessary baseline rather than a sufficient validation criterion, [Yi- what does that mean?] and next evaluate each model using process probes whose observables depend on the shape of the potential-energy surface along physically motivated paths. [Yi- such as?]
% REV-FORMALIZATION-2026-09-12 END.

\begin{figure}[ht]
    \includegraphics[width=0.8\textwidth]{Fig/NeurIPS2025-AI4MAT-fig2-v2.pdf}
    \centering
    \caption{\textbf{Comprehensive Benchmarking of MLFF Performance in Molecular Dynamics Simulations.} (a) A schematic of the unified evaluation protocol. Four candidate models (a zero-shot foundation model, a bespoke model trained from scratch, and two fine-tuned variants) were used to drive a 200~ps MD simulation. The resulting trajectories are then subjected to a uniform set of post-processing analyses to evaluate key physical properties. (b) Thermodynamic stability, demonstrated by the evolution of temperature and pressure over the simulation, remain acceptably stable around their target values for all models. (c) Thermal transport properties, showing the Heat Flux Autocorrelation Function (HFACF) and its running integral to compute thermal conductivity ($\kappa$).  (d) The Mean Squared Displacement (MSD), used to assess atomic mobility and calculate the diffusion coefficient. (e) The Velocity Autocorrelation Function (VACF), which describes the system's underlying dynamics.}
    \label{fig:ensemble}
\end{figure}


All MACE models (Foundation, Scratch, FT-600K, FT-Multi\_T) successfully reproduce the equilibrium structure of Cr-doped Sb$_2$Te$_3$, with RDF analysis showing excellent agreement with the AIMD reference (``ground truth'') across all atomic pair correlations (Fig.~S2). This Figure includes a structural schematic of the local environment around Cr and full radial distribution functions (RDFs) for all models, confirming that each training strategy preserves the characteristic nearest-neighbor and second-shell coordination features at 600~K. While all models maintain stable thermodynamic ensembles at 600~K (Fig.~\ref{fig:ensemble}c), the zero-shot foundation model shows a persistent pressure offset, likely originating from its exposure to predominantly 0~K equilibrium structures rather than configurations extant at finite temperatures.

\reviewermap{REV-C9}{pending-cluster}{Long-run MD transport values must be replaced by the corrected 100000-step workflow before final submission.}

Despite similar performance on local structural (RDF) and short-time dynamic (VACF) properties (Fig.~\ref{fig:ensemble}e), the models diverge in their preliminary predictions of longer timescale transport phenomena. The original FT--600K transport calculation used a non-comparable trajectory protocol [Yi- how do you know that about the trajectory?] and is therefore quarantined [Yi- what does ``quarantined'' mean in this context?] rather than used for final model ranking. Because diffusion and thermal transport coefficients are sensitive to sampling length and trajectory health, the final analysis will replace these values with the corrected matched long-run MD workflow [PLACEHOLDER: corrected 100000-step MD diffusion and thermal conductivity values]. If the audited ranking is unchanged, we will retain a qualified interpretation that high-temperature training exposure can flatten the effective PES; if the ranking changes, the text and figure will be regenerated from the audited workflow. [Yi- I need you to explain to me what this last sentence means]

The current thermal conductivity curves are therefore treated as a protocol-audit target rather than a final transport conclusion. They still motivate the methodological point that validating MLFFs solely on static structures and short-time dynamics is insufficient: collective transport properties require matched simulation protocols, uncertainty estimates, and trajectory health checks before they can support model-ranking claims.




\subsection{Non-Equilibrium Diffusion Pathways: The Critical Test of Generalization}

% REV-FORMALIZATION-2026-09-12 BEGIN: replace broad generalization language with explicit path-operator framing.
% Original paragraph began: "A rigorous test of the generalizability..."

The migration path benchmark evaluates a different measurement operator from the equilibrium RMSE. For each selected end-point pair $\tau=(a_s,a_t,c)$, the DFT NEB path defines a process-conditioned sequence of configurations $\Gamma_{\mathrm D,\tau}^{*}$ that includes the metastable endpoints, intermediate distorted images, and the transition-state region. A candidate MLFF is then tested either by fixed-path scoring on this DFT geometry or by self-relaxed NEB optimization on its own PES, as defined in Sec.~\ref{sec:neb_operators_formal}. These two operators answer different questions: whether the model assigns correct energies to the reference path, and whether the model can stably discover a plausible path without being given the DFT images.

% REV-FORMALIZATION-2026-09-12 END.

\begin{figure}[ht]
    \includegraphics[width=0.8\textwidth]{Fig/NeurIPS2025-AI4MAT-figure4-v2.pdf}
    \centering
    \caption{\textbf{Comparative analysis of MLFF training strategies for predicting atomic migration barriers.} (a) Schematic illustration of the simulated Cr atom migration pathway between two stable sites within the Sb$_2$Te$_3$ bilayer. (b) Bar plot showing the migration energy prediction errors for various models, spanning bespoke training to advanced active- learning approaches. (c--f) Comparison of the minimum energy pathway (MEP) profiles for the Cr migration process. The solid black line denotes the ground-truth DFT reference data, while dashed lines represent predictions from MACE models trained with different strategies. 
}
    \label{fig:barrier}
\end{figure}


\subsubsection{Local Migration Events}

% REV-FORMALIZATION-2026-09-12 BEGIN: shorten local path setup now that fixed-path operator is formally defined.
% Original paragraph began: "To perform a direct and controlled evaluation..."

For the fixed-path component of the benchmark, we used the DFT-relaxed NEB images as a common geometric reference and evaluated each MLFF by single-point energies on those images. This corresponds to $\Delta E_{\mathrm{fixed}}^{\ddagger}(M,\tau)$ in Sec.~\ref{sec:neb_operators_formal}: the model is not allowed to move the path, so the resulting error isolates energetic fidelity along a DFT-supplied trajectory.
% REV-FORMALIZATION-2026-09-12 END.

\reviewermap{REV-A1}{partial}{Fig. 3d is fixed-geometry DFT-path scoring; current 1-4 DFT reference candidate is 0.336050 eV; Foundation fixed-path error is +0.688246 eV.}
% REV-FORMALIZATION-2026-09-12 BEGIN: report candidate DFT barrier and Foundation fixed-path error under the formal operator.
% Original paragraph began: "Here the DFT NEB path is used..."

For the in-gap pathway, labeled 1--4 in the audited workflow, the current DFT candidate barrier is 0.336~eV. This value will be retained only after the restarted QE NEB run satisfies the stated convergence criterion; otherwise the table and plot will be regenerated from the converged restart. [Yi- Please explain what this sentence means too] The MACE Foundation fixed-path barrier is 1.024~eV, [Yi- do we have to quote these energy barriers to 3 d.p.? They aren't that accurate] giving an audited fixed-path error of +0.688~eV relative to the current DFT candidate. Thus, although the foundation model gives a smooth qualitative profile, it is not quantitatively accurate on this reference path.

% REV-FORMALIZATION-2026-09-12 END.
% REV-FORMALIZATION-2026-09-12 BEGIN: compress baseline interpretation to avoid duplicating the fixed-path operator values.
% Original paragraphs began: "The two baseline strategies..." and "Conversely, the MACE Foundation model..."



The two baseline strategies illustrate different fixed-path failure modes. The MACE Scratch model strongly overestimates the barrier, consistent with insufficient sampling of high-energy, low-probability transition-state configurations in the training data. The MACE Foundation model produces a smoother qualitative profile but remains quantitatively miscalibrated on the DFT-supplied path, consistent with the possibility that broad pretraining captures generic PES smoothness while missing system-specific curvature along a rare-event pathway.~\citep{deng2025softening} 
These results motivate fine-tuning, while the formal separation of fixed-path scoring and self-relaxed NEB optimization remains essential for interpreting the comparison. [Yi- please explain this sentence]
% REV-FORMALIZATION-2026-09-12 END.


% \paragraph{Further Testing Different Trajectories}

To probe the capability of generalizing the different MLFF training strategies, we evaluated their performance on two distinct Cr migration pathways (Fig.~\ref{fig:barrier}a) with fundamentally different characteristics. The first pathway, \emph{in-gap diffusion}, involves Cr migration within the vdW gap between Sb$_2$Te$_3$ quintuple layers, with low energy barriers and configurations similar to the training data, thereby testing interpolative accuracy. [Yi- more explanation please]

In contrast, the second pathway, \emph{deep penetration}, requires the Cr atom to move vertically from the vdW gap and penetrate directly into the interior of a QL, ultimately reaching a deeply buried site within the layer. This migration path is associated with significantly higher energy barriers and accesses high-distortion configurations that are not well-represented in the training data, making it an out-of-distribution (OOD) scenario with potentially more metastable configurations along the pathway. Testing both pathways allows us to evaluate model accuracy on familiar configurations and extrapolative power in challenging regions of the configuration space.


\subparagraph{In-Gap Diffusion: A Test of Interpolative Accuracy}

For the in-gap diffusion pathway (Fig.~\ref{fig:barrier}c--d), in which the Cr atom moves within the vdW gap or just shallowly penetrates the interface, the atomic environments along the minimum energy path (MEP) are expected to be closer to the near-equilibrium states sampled during the AIMD simulations than the deep-penetration path. Under the grouped-split, multi-seed re-evaluation (seeds 123/234/345), the absolute fixed-path barrier error on pathway 1--4 relative to the current DFT candidate (0.336~eV) is $0.474 \pm 0.121$~eV for FT--600K fine-tuned from the OMAT-FT checkpoint, $1.204 \pm 0.270$~eV for Scratch, and $0.678 \pm 0.195$~eV for Scratch-5\%; a single-seed MACE-MPA-0 fine-tuned comparator gives $0.237$~eV but is excluded from these seed-uncertainty statistics. FT--600K therefore remains the smallest-error fine-tuned model on this pathway, consistent with the single-model result reported previously, but the grouped, multi-seed re-evaluation does not reproduce the earlier point estimate of a 0.16~eV barrier error: the audited error under the present protocol is larger and has seed-to-seed spread. We treat the ranking among FT--600K, Scratch, and Scratch-5\% on this single pathway as provisional evidence for task-relevant fine-tuning rather than a precise barrier-error claim, and we do not extend this pathway-1--4 comparison to the other audited pathways or to FT--MultiT, which was not retrained under the grouped protocol in this round.

To place the specialist models on the same fixed-path scale as current foundation checkpoints, we additionally evaluated ten single-checkpoint foundation models using the identical five-image pathway-1--4 inputs. The four representative checkpoints (MP-0b3, MPA0, OMAT-0 medium, and MH-1) predicted barriers of 1.503, 1.791, 1.602, and 1.841~eV, respectively, compared with the provisional DFT value of 0.336050~eV. Thus, all four foundation representatives overestimate the barrier by 1.17--1.50~eV; MPA0 is not the kinetic-accuracy winner in this completed comparison. The individual profiles and the family-aggregated profiles are provided in the Supplementary Information (Fig.~S\ref{fig:foundation_profiles}), while the full preflight, raw outputs, and fixed-path diagnostics are archived with the benchmark release. This comparison is deliberately reported as single-checkpoint evidence: foundation models do not provide seed-based uncertainty, and the declared overlap audit does not test unknown foundation-pretraining corpora.

\subparagraph{Deep Penetration: A Test of Extrapolative Robustness}

The ``deep penetration'' pathway provides a much more stringent test of model robustness. This path involves significant lattice distortion as the Cr atom pushes its way through the covalently bonded quintuple layer (QL), creating atomic environments that are far from those seen in the training data (Fig.~\ref{fig:barrier}a).

% REV-FORMALIZATION-2026-09-12 BEGIN: soften deep-penetration claim and tie it to mu_tau/OOD evidence level.
% Original paragraph began: "While the fine-tuned and foundation models remain accurate..."
The deep-penetration path samples a different region of $\mu_\tau$ from the in-gap path: its transition-state images involve larger lattice distortion and weaker overlap with the AIMD training distribution. We therefore use it as a probe of extrapolative robustness rather than as a direct extension of the near-equilibrium RMSE test. Until the restarted QE references and multi-seed model rankings are closed, [Yi- needs explanation] the conservative conclusion is that near-equilibrium accuracy does not determine rare-event barrier accuracy.
% REV-FORMALIZATION-2026-09-12 END.

\reviewermap{REV-A3}{pending-cluster}{Grouped, multi-seed retraining of Scratch and FT-600K is complete (Sec.~3.3.3 in-gap results); deep-penetration pathway-ID classification and restarted QE reference remain pending before ranking models on this OOD probe.}

The MACE Scratch model, which performed least well of the four models on the in-gap pathway, yields the lowest barrier error for this out-of-distribution task (Fig.~\ref{fig:barrier}e, f). We do not interpret this isolated result as evidence that the scratch model is physically superior. Rather, because the deep-penetration path samples highly strained configurations and the current DFT NEB trajectories are still being consolidated, this comparison is treated as a question of seed- and path-stability. The completed analysis will report the multi-seed barrier distribution and the converged QE reference values for these paths [PLACEHOLDER: multiseed deep-penetration barrier errors and convergence status]. Until those checks are complete, the result supports the narrower conclusion that equilibrium accuracy does not reliably predict rare-event barrier accuracy.

\paragraph{The Critical Role of Task-Specific Fine-Tuning}
\reviewermap{REV-A5}{partial}{Scratch-5\% control (3 seeds) is trained, RMSE-evaluated, and evaluated on the in-gap pathway (Sec.~3.3.3); multi-pathway comparison remains pending.}

% REV-FORMALIZATION-2026-09-12 BEGIN: preserve FT-600K result while framing it as case-study process-fidelity evidence.
% Original paragraph began: "The dramatic improvement seen in the fine-tuned models..."

The improvement seen in the fine-tuned models, especially for the in-gap diffusion task (Fig.~\ref{fig:barrier}b), shows that adaptation can improve process fidelity when the fine-tuning data overlap with the task-conditioned pathway distribution. The MACE FT--600K model, fine-tuned on data generated at the target temperature of the migration event, gives the smallest reported in-gap barrier error in the current analysis. This supports the interpretation that task-relevant fine-tuning can recalibrate the local PES along a selected migration probe, while leaving open whether the same adaptation improves all non-equilibrium processes.
% REV-FORMALIZATION-2026-09-12 END.

% REV-FORMALIZATION-2026-09-12 BEGIN: replace broad "more data" claim with path-aware task-relevance claim and acceptance criteria.
% Original paragraph began: "Intriguingly, the MACE FT--Multi-T model..."

Within this case study, task relevance appears to be more predictive than data volume alone: the 600~K fine-tuned model is closer to the in-gap migration probe than the broader multi-temperature fine-tuned model. This interpretation depends on path-aware train--test proximity and the Scratch-5\% control, which are included in the acceptance criteria in Sec.~\ref{sec:acceptance_formal}. The grouped-split, nearest-neighbor SOAP/RMSD audit described above finds no in-gap NEB-image overlap with either the full grouped training set or the Scratch-5\% subset. The Scratch-5\% control (same $\sim$1,000-configuration grouped FT-600K training subset, random initialization, seeds 123/234/345) gives a pathway-1--4 absolute barrier error of $0.678 \pm 0.195$~eV, intermediate between FT--600K-OMATFT ($0.474 \pm 0.121$~eV) and full Scratch ($1.204 \pm 0.270$~eV) on this single pathway; given the small sample (one pathway, three seeds, $\sim$90 training frames), we treat this as consistent with, rather than definitive proof of, a foundation-pretraining benefit beyond data volume alone. The corresponding comparison against FT--MultiT could not be updated in this round because FT--MultiT was not retrained under the grouped protocol.
% REV-FORMALIZATION-2026-09-12 END.

% REV-FORMALIZATION-2026-09-12 BEGIN: soften "prove" language and connect data selection to benchmark-guided acquisition.
% Original paragraph began: "The success of targeted fine-tuning points the way forward..."
The success of targeted fine-tuning points the way forward, while the divergent results of the 600~K and Multi-T models indicate that the data selection process is non-trivial. This motivates benchmark-guided data acquisition, including uncertainty-guided active learning, to identify configurations that improve model robustness and accuracy in a data-efficient manner.
% REV-FORMALIZATION-2026-09-12 END.


\subsubsection{NEB Stability as an Additional Robustness Criterion}
\reviewermap{REV-A1}{partial}{Define Fig. 3d and SI Fig. S3 as different NEB operators and report them side-by-side.}
% REV-FORMALIZATION-2026-09-12 BEGIN: align Fig. 3d/S3 explanation with fixed-path and self-relaxed operators.
% Original paragraph began: "Beyond evaluating the models on pre-computed DFT NEB geometries..."

We also evaluated the separate self-relaxed NEB operator $\Delta E_{\mathrm{self}}^{\ddagger}(M,\tau)$, in which the MLFF supplies the forces used to optimize the band. The 0.41~eV energy barrier value reported in Fig.~S3 is the MACE Foundation model's self-consistent barrier on its own optimized path. It is, therefore, not a second DFT reference and not a contradiction of the +0.688~eV fixed-path error in Fig.~\ref{fig:barrier}d. In contrast, the scratch, FT--600K, and FT--MultiT models exhibit force growth and early termination in this self-relaxed task, [Yi- we need to discuss this sentence. what is ``force growth'' and early termination?] making NEB stability itself an additional robustness criterion. Any low-barrier MLFF-relaxed path should be interpreted physically only after DFT single-point or DFT-restarted validation.
% REV-FORMALIZATION-2026-09-12 END.

\begin{figure}[ht]
    \includegraphics[width=0.6\textwidth]{Fig/NeurIPS2025-AI4MAT-figure5.pdf}
    \centering
    \caption{\textbf{Evaluation of MACE models on collective lattice displacement: interlayer sliding in Sb$_2$Te$_3$.} (a) Schematic illustration of the bilayer sliding process in Sb$_2$Te$_3$, where the top layer (yellow/orange atoms) slides relative to the bottom layer along the crystallographic direction. The purple sphere indicates the position of a Cr dopant when present. (b) Energy barriers for interlayer sliding in pristine Sb$_2$Te$_3$ as calculated by DFT (black, ground truth) and various MACE models. The reaction coordinate represents the normalized sliding distance from the initial to final configuration. (c) Energy barriers for the same sliding process in Cr-doped Sb$_2$Te$_3$.}
    \label{fig:sliding_results}
\end{figure}



\subsection{Interlayer Sliding: Probing Robustness to Non-Local, Collective Displacements}

Having tested the models' response to local, high-distortion events, we evaluated their robustness to a different class of out-of-distribution challenge: large-scale, collective lattice displacements. To this end, we simulated the shearing of one Sb$_2$Te$_3$ layer relative to another, a process governed by the subtle corrugations of the interlayer vdW potential energy surface (Fig.~\ref{fig:sliding_results}a). This task is particularly challenging for MLFFs as the local atomic environment of any given atom changes only minimally, while the global configuration undergoes a significant transformation that crosses periodic boundaries. The configurations along this sliding path were not present in the training data.

We first examined the case of pure, undoped Sb$_2$Te$_3$ (Fig.~\ref{fig:sliding_results}b). The results reveal a clear trade-off between the models. The MACE Foundation model provides the most reasonable estimate of the energy barrier's shape and magnitude, suggesting its vast pre-training on bulk crystals has endowed it with a better implicit understanding of such periodic, mechanical deformations. However, it fails to maintain translational symmetry, incorrectly predicting the final state to be higher in energy than the identical initial state. This energy drift is a clear artifact, indicating a failure to perfectly respect the periodic nature of the simulation cell under large displacements. The MACE Scratch model exhibits a similar version of this artifact.

In contrast, the fine-tuned models (MACE FT--600K and Multi-T) significantly underestimate the energy barrier. This suggests a testable hypothesis: fine-tuning that improves local chemistry around the Cr dopant may reduce fidelity for weaker, longer-range interlayer physics inherited from the foundation model. The present data support this interpretation as a case-study observation rather than as a universal statement about fine-tuning.

This trend is also observed in the Cr-doped system (Fig.~\ref{fig:sliding_results}c), suggesting that it is not solely an artifact of the pristine sliding setup. The failure of all models to perfectly capture both the barrier height and the end-point periodicity highlights a limitation relevant to local-descriptor-based MLFFs. Phenomena like shearing, stacking faults, and dislocation glide are inherently non-local. While the models can describe local bonding and coordination, they may struggle to enforce long-range physical constraints that extend beyond their cut-off radius. This further underscores the need for careful validation when using standard MLFFs to study mechanical properties and points towards future work in developing training sets that explicitly include these collective deformation modes or exploring architectures designed to capture long-range physics more effectively.

\subsection{Representation Learning Analysis}

To diagnose possible origins of the observed performance differences, we conducted a multi-faceted analysis of the learned representations. We projected high-dimensional atomic environment descriptors from each model's 600~K MD trajectory onto a two-dimensional space using two complementary visualization techniques: $t$-distributed stochastic neighbor-embedding ($t$-SNE)~\citep{maaten2008visualizing} to assess representational dissimilarity, and Potential of Heat-diffusion for Affinity-based Trajectory Embedding (PHATE)~\citep{moon2019visualizing} to visualize trajectory geometry (Fig.~\ref{fig:tsne-phate}).

The $t$-SNE projection (Fig.~\ref{fig:tsne-phate}a) suggests that models trained with different strategies produce qualitatively distinct encodings. The MACE Foundation and MACE Scratch models occupy separated regions in the projection, while the two fine-tuned models occupy an intermediate region. Because low-dimensional embeddings can amplify or suppress apparent separation, we use these plots as qualitative diagnostics and report quantitative checks in the original descriptor space. In the audited analysis, the original 6191-dimensional Euclidean silhouette score is 0.333, while z-scoring the original features gives a silhouette score of -0.019; the projected baselines are $0.397\pm0.027$ for $t$-SNE and 0.473 for PHATE. These values support representational differences but they do not, by themselves, prove a unique mechanistic cause.

PHATE provides a complementary visualization of trajectory geometry, but we do not treat the projection as a direct mechanistic explanation. The PHATE projection (Fig.~\ref{fig:tsne-phate}b) maps the continuous time-evolution of the system to a low-dimensional trajectory manifold whose geometry is compared with the potential energy in Fig.~\ref{fig:tsne-phate}c. Along the primary axis of this manifold (PHATE 1), the specialist models (Scratch and FT) occupy similar regions, while differences along PHATE 2 separate the scratch-trained model from the fine-tuned models.

This separation is consistent with, but does not prove, a difference in how the model families represent configurations sampled during the dynamics at 600~K. The scratch-trained model's isolated region may reflect a narrower or less constrained representation, whereas the fine-tuned models' intermediate locations are consistent with adaptation from the pre-trained foundation model. Direct causal claims require perturbation or intervention tests beyond these projections.

\reviewermap{REV-A6/REV-C11/REV-C12}{done-local}{Claims softened to consistent-with; original 6191-d euclidean silhouette is 0.333; zero feature stub dropped.}

This geometric difference in representation is consistent with the observed diffusion and NEB behavior, but we do not treat the embedding as a direct mechanistic proof. Diffusion is a rare-event process that requires the model to generalize to out-of-distribution transition states. In this context, the scratch model's separated representation is consistent with a less constrained PES outside its training domain, while the fine-tuned models' intermediate representations are consistent with adaptation of a general physical prior to the target system. We therefore interpret the latent-space analysis as a diagnostic for model families that should be read together with the task-level barrier, transport, and stability benchmarks.

\begin{figure}[ht]
    \includegraphics[width=\textwidth]{Fig/NeurIPS2025-AI4MAT-figure6-phate.pdf}
    \centering
    \caption{\textbf{Representation diagnostics from trajectory embeddings and original-descriptor checks.} Projections of atomic environment descriptors from 600~K MD trajectories. 
    \textbf{(a)}~A $t$-SNE projection shows that Foundation (red) and Scratch (orange) models produce highly separated representations, while fine-tuned models (blue) are intermediate. 
    \textbf{(b)}~PHATE projection visualizes the system's continuous trajectory manifold. The separation between the scratch-model and fine-tuned-model projections along PHATE 2 is interpreted as a qualitative diagnostic rather than direct mechanistic proof. 
    \textbf{(c)}~The same PHATE embedding, colored by potential energy, compares projection geometry with model energy values. The foundation model's large energy offset indicates its nature as an uncalibrated prior.
\reviewermap{REV-C11}{done-local}{Silhouette values are recomputed in original descriptor space; projected values are retained only as a baseline.}
    \textbf{(d)}~Silhouette scores quantify separability in the original descriptor space and in low-dimensional embeddings. The projected values are used as visualization diagnostics, while the original-space values provide the primary quantitative check on representational dissimilarity.}
    \label{fig:tsne-phate}
\end{figure}


%%%%%%%%%% new parts for SHAP %%%%%%%%%%%%%


\subsection{Dissecting Learned Representations with Feature-Level Interpretability}

While $t$-SNE and PHATE suggest that distinct training strategies produce separable descriptor patterns, these projections do not establish the physical principles learned by each model. To add a feature-level diagnostic, we use SHapley Additive exPlanations (SHAP)~\cite{lundberg_unified_2017} on surrogate models of MACE's absolute energy error. This analysis identifies SOAP descriptors to which the surrogate error predictors are most sensitive; it does not directly explain the internal message-passing mechanism of MACE.

\reviewermap{REV-A6/REV-C13}{done-local}{SHAP explains surrogate error predictions, not MACE internals directly; current 5-fold CV R2 values are 0.9819, 0.8756, and 0.9730. The X-FORCE three-model lineage is recorded in the release audit.}

Our methodology is designed to identify descriptors correlated with predictive error. For each MACE model, we trained a Gradient-Boosting Regressor to predict absolute total-energy error from 390 SOAP and 17 structural descriptors extracted from the X-FORCE 150-frame trajectory; frame identity is joined explicitly through \texttt{image\_index}.~\cite{bartok2012representing} This approach interprets a simpler surrogate of the error landscape, not the internal message-passing mechanism of MACE itself. The audited 5-fold cross-validation $R^2$ values are 0.9819, 0.8756, and 0.9730 for FT--600K, FT--Multi--T, and Scratch, respectively. We then applied TreeSHAP to these surrogate models, assigning feature-importance values to SOAP descriptors that are most correlated with prediction errors.


Our analysis shows a divergence in the surrogate feature-importance landscapes associated with each training strategy. A global overview of the top 50 influential descriptors across models (Fig.~\ref{fig:shap}b) shows that the fine-tuned multi-temperature (FT--Multi-T) surrogate concentrates importance more strongly, whereas the scratch-model surrogate distributes importance more diffusely. In the locked revised ranking, the largest FT--Multi-T contribution is \texttt{Cr-Sb\_n43\_l0}; the leading Scratch contributions are dominated by Cr--Te descriptors. These are descriptive surrogate sensitivities, not evidence that the descriptors are causal or mechanistically internalized by MACE.
% This disagreement maps directly onto the differences in model performance: the fine-tuned models suppress these unstable, non-consensus descriptors, whereas the scratch-trained model relies on them—revealing that its poor performance originates from assigning importance to features that do not consistently encode physically meaningful interactions.

This disparity is quantified in Fig.~\ref{fig:shap}a, which compares the feature-importance profiles of the fine-tuned and scratch models. The FT--Multi-T feature is consistent with sensitivity to Cr--Sb local coordination in the doped lattice, but a direct physical assignment requires perturbation tests rather than SHAP rankings alone. The scratch model's broader surrogate-importance profile is consistent with a less concentrated error predictor; it does not by itself establish a physical explanation for the model's NEB behavior.

% A broader, physics-based categorization of the features (Fig.~\ref{fig:shap}c) reinforces this finding. The FT--Multi-T model consistently assigns the highest importance to features characterized by long-range and high angular complexity. In contrast, the scratch model distributes its attention more evenly and ineffectively, particularly across simpler short-range and spherical interactions. This demonstrates that the fine-tuning process does not merely adjust parameters; it fundamentally re-weights the model's attention, distilling knowledge from the pre-trained foundation model to prioritize physically dominant interactions. This interpretability framework provides a quantitative, feature-level explanation for the model's learned chemical intuition and offers a powerful diagnostic tool for optimizing MLFF training protocols for specialized chemical systems.

%%%


\begin{figure}[h!]
    \centering
    \includegraphics[width=0.95\linewidth]{revision1/figures/fig6_shap_revised.pdf}
    \caption{\textbf{SHAP-based Dissection of Feature Importance in MACE Models.}
    \textbf{(a)}~Mean absolute SHAP values for the top 50 SOAP features, comparing the best-performing fine-tuned model (FT--Multi-T) and the worst-performing scratch-trained model. The corrected audit identifies \texttt{Cr-Sb\_n43\_l0} as the largest mean absolute SHAP feature in the FT--Multi-T surrogate.
    % \yi{explicitly say the difference between the orbital }
    \textbf{(b)}~Heat map of surrogate feature importance for the top 50 SOAP features across all training strategies. The FT--Multi-T surrogate concentrates importance on fewer features, while the scratch-model surrogate distributes importance more diffusely.
    % \textbf{(c)}~Aggregated SHAP importance across six physically-meaningful SOAP feature categories. The fine-tuned models consistently prioritize features with long-range and high angular complexity, indicating a more sophisticated learned representation of the underlying chemical physics.
    }
    \label{fig:shap}
\end{figure}

\section{Conclusions}

% REV-FORMALIZATION-2026-09-12 BEGIN: conclusion now cites the formal measurement-operator framework.
% Original paragraph began: "We benchmarked specialist (from-scratch)..."

We benchmarked specialist and foundation-based MLFFs for Cr-doped \ce{Sb2Te3} using the measurement-operator framework defined in Sec.~\ref{sec:benchmark_formalization}. This formalism separates equilibrium point-wise accuracy from task-conditioned process fidelity, and it distinguishes fixed-path energetic scoring from self-relaxed path optimization. Under this framework, diffusion trajectories and structured pathways show that models with similar equilibrium metrics can diverge sharply under kinetically relevant distortions.
% REV-FORMALIZATION-2026-09-12 END.

\reviewermap{REV-A4}{done-local}{Broader principles are now framed as case-study hypotheses requiring cross-system validation.}

Our findings support several broader hypotheses. Robust benchmarking should span both in-distribution and out-of-distribution tasks, as accuracy near equilibrium does not guarantee kinetic fidelity. Fine-tuning enhances targeted performance but can reduce global physical generality, emphasizing the need for carefully designed adaptation strategies. [Yi- can we make this sentence more easy to understand] Latent space analysis shows that different training paradigms can produce qualitatively distinct representations, helping to diagnose why models with similar traditional metrics may behave differently in downstream tasks. Notably, the deep penetration paths remain the most stringent OOD regime in this benchmark and are being finalized through restarted QE NEB calculations before final numerical ranking.

\reviewermap{REV-A4}{done-local}{Conclusion should keep framework generality as future validation rather than a final claim.}
Together, these results indicate that adding ``more data'' is not necessarily sufficient; task-relevant, non-equilibrium sampling is often the decisive ingredient for meaningful improvement. Migration-inspired probes, coupled with latent space diagnostics and auditable data provenance, provide a practical route to identify which additional configurations are needed and where models are likely to fail. In this case study, the framework reframes the specialist--generalist debate: the goal is not to choose \textit{between} model classes, but to exploit the complementary strengths they offer while validating them against the specific non-equilibrium tasks they are expected to perform.

\section*{Acknowledgments}

\begin{itemize}
  \item \textbf{Funding} \\
This work was supported by the United States Department of Defense-funded Center of Excellence for Advanced Electro-photonics with 2D materials—Morgan State University, under Grant No. W911NF2120213. The authors thank Drs. Frank Gardea and Owen Vail, the cooperative agreement managers of the Center, for their interest and constructive comments. They also thank Dr. Ramesh Budhani for his guidance on the project.
YC thanks Johns Hopkins University for their support in AY 24–25. Computational resource support was provided by the petascale Hopkins facility, Advanced Research Computing at Hopkins (ARCH) (rockfish.jhu.edu), supported by National Science Foundation award OAC 1920103. Partial funding for ARCH’s infrastructure was originally provided by the State of Maryland.

  \item \textbf{Conflict of interest/Competing interests} \\
  The authors declare that they have no known competing financial interests or personal relationships that could have appeared to influence the work reported in this paper.

  \item \textbf{Ethics approval and consent to participate} \\
  Not applicable. This study does not involve human participants or animal subjects.

  \item \textbf{Consent for publication} \\
  Not applicable.

  \item \textbf{Data availability} \\
  The supporting dataset, including training manifests, NEB path metadata, model-predicted barriers, and post-processing tables, will be released through MigrationBench on Hugging Face at \url{https://huggingface.co/datasets/alinacao2000/MigrationBench} [PLACEHOLDER: final dataset DOI or commit hash]. The manifest records the source trajectory, split assignment, random seed, model provenance, DFT/MLFF task type, convergence status, and checksum for each derived artifact. NEB paths are evaluation/reference data and were not merged into the fine-tuning datasets. The frame-level force audit covers 9,163 released training frames and found no unclassified or NEB-derived all-zero-force frame; the nine exact zero-force records are explicitly labelled isolated-atom $E_0$ references (release commit \texttt{f5c6d19}).

  \item \textbf{Materials availability} \\
  Not applicable. This work is entirely computational and does not involve the synthesis or distribution of physical materials.

  \item \textbf{Code availability} \\
\reviewermap{REV-C8/REV-C13}{done-local}{The code release includes the fail-closed zero-force exporter fix, frame-level provenance audit, and X-FORCE-derived SHAP release.}
  The scripts used for data processing, DFT/MD post-processing, and plotting are available on GitHub \url{https://github.com/yicao-elina/MigrationBench.git}. The Sec.~3.5/Fig.~6 SHAP chain is released as \texttt{scripts/shap\_analysis.py}, \texttt{scripts/shap\_feature\_labels.py}, and \texttt{scripts/plot\_shap\_fig6.py}, using the compact three-model X-FORCE inputs in \texttt{data/raw\_compact/revision1/shap/}. The source scope and MD5s are recorded in \texttt{docs/audit/shap\_source\_lineage.md}; the four-model, 2-concentration, and Dual-X ID/OOD/NEB workflows are not inputs to Fig.~6.

  \item \textbf{Author contribution} \\
    \textbf{Y. C.}: conceptualization, methodology, data curation, DFT and MD simulations, analysis, visualization, and manuscript-writing. \textbf{P. M.}: literature review, conceptualization, and manuscript- writing. \textbf{P. C.}: conceptualization, project administration, supervision, funding acquisition, and manuscript-writing.
\end{itemize}


\bibliography{sn-bibliography}
% \bibliographystyle{sn-nature}



% \newpage
% \appendix


\end{document}
